\documentclass[UTF8,a4paper,11pt]{ctexart}

\usepackage{amsmath}
\usepackage{amssymb}
\usepackage{booktabs}
\usepackage{enumitem}
\usepackage{geometry}
\usepackage{graphicx}
\usepackage{hyperref}
\usepackage{listings}
\usepackage{svg}
\usepackage{xcolor}

\geometry{margin=2.4cm}
\hypersetup{
  colorlinks=true,
  linkcolor=blue,
  urlcolor=blue,
  citecolor=blue
}

\definecolor{codebg}{RGB}{248,248,248}
\definecolor{codekw}{RGB}{0,92,175}
\definecolor{codecomment}{RGB}{80,120,80}

\lstdefinelanguage{mychem}{
  morekeywords={Ph,Im,Ind,Pyr,Ox5,Me,Ac,NH2,OH,SH,OMe,OPh,C,N,O,S,P,F,Cl,Br,I},
  sensitive=true,
  morecomment=[l]{\#}
}

\lstset{
  basicstyle=\ttfamily\small,
  backgroundcolor=\color{codebg},
  keywordstyle=\color{codekw},
  commentstyle=\color{codecomment},
  breaklines=true,
  columns=fullflexible,
  keepspaces=true,
  frame=single,
  framerule=0.2pt,
  rulecolor=\color{black!25},
  showstringspaces=false
}

\svgsetup{inkscapelatex=false}

\title{MarkChem：面向大模型化学回答的可编译分子标记语言}
\author{何梓源}
\date{\today}

\begin{document}

\maketitle

\begin{abstract}
本文报告一个面向化学结构表达与网页渲染的小型编译系统。系统定义了一种文本 DSL，将大模型或人工给出的分子片段表达式编译为显式的分子连接图 IR，并进一步渲染为 SVG 键线式。本文目前重点描述技术模块，包括语言语法、语义设计、AST 到 IR 的翻译、前后端划分以及 C++ 实现。研究动机、实验结果和讨论将在后续版本中补充。
\end{abstract}

\section{动机}

\begin{itemize}[leftmargin=2em]
  \item 内嵌LaTeX公式的Markdown是大模型输出的事实标准，其对数学/物理公式有着很好的支持，不仅可以由AI生成，人类经过简单学习也可以看懂并编辑。
  \item 目前现的对化学领域常见的结构表达方式大致有以下几种
        \begin{table}[h]
        \centering
        \begin{tabular}{lll}
        \toprule
        例子 & 优势 & 劣势 \\
        \midrule
        \verb|SVG| & \verb|直观、生成成本较低| & 需要模型多模态能力、编辑困难、丢失了高级语义 \\
        \verb|SMILES| & \verb|化学领域数据库常用标准| & 人类难懂、编辑困难 \\
        \verb|InChI| & \verb|具有唯一性，索引方便| & 语法限制性强，图到文转化困难 \\
        \verb|ToolCall| & \verb|直接接入专业软件、正确稳定| &  现有支持不多，不利于构建完整生态\\
        \bottomrule
        \end{tabular}
        \end{table}
  
  \item 于是我希望提出一种AI可生成、人类可手写，检查，局部编辑、适合网页内嵌的DSL。
\end{itemize}

\section{系统概览}

系统将分子结构表达建模为一次小型编译过程：
\[
  \text{my-Chem source}
  \rightarrow \text{Token stream}
  \rightarrow \text{AST}
  \rightarrow \text{IR graph}
  \rightarrow \text{SVG}.
\]

前端负责识别和解析源语言，后端负责语义翻译与渲染。整体实现由一个 C++ 编译核心库和若干入口程序组成：
\begin{itemize}[leftmargin=2em]
  \item \textbf{词法分析器}：将输入字符串切分为括号、分隔符、数字、关键字等 token。
  \item \textbf{递归下降 Parser}：根据文法构造 AST，保留 CPO、接口、核心片段和取代基结构。
  \item \textbf{语义翻译器}：把 AST 翻译为无序分子连接图 IR，并完成接口合法性、端基、稠合、原子替换等语义检查。
  \item \textbf{SVG 渲染后端}：优先将 IR 降低为 RDKit 分子对象并调用 RDKit 2D depiction；当 RDKit 不可用时使用内置 SVG fallback。
  \item \textbf{浏览器插件与本地服务}：网页中的 \verb|<my-Chem>...</my-Chem>| 内容由插件提取，发送给本地 C++ 服务，返回 SVG 后替换原标签。
\end{itemize}

\section{源语言设计}

\subsection{CPO 片段}

源语言的基本单位是 CPO（chemical phrase object）。一个完整输入由一个或多个 CPO 串联组成，每个 CPO 以分号结束：

\begin{lstlisting}[language=mychem]
CPD       := CPO+
CPO       := INTERFACE CORE "," SUBS INTERFACE ";"
INTERFACE := "(" POSES ")" | "{" POSES "}"
POSES     := digit ("," digit)*
\end{lstlisting}

一个 CPO 包含四个部分：左接口、核心片段、取代基列表、右接口。例如：

\begin{lstlisting}[language=mychem]
(0)Ph,1-OH,(0);
(0)NH2CC,4-Me,5-=O,5-OH,(0);
(0)6R,{4,5};{1,2}6R,(0);
\end{lstlisting}

接口中的编号均为当前 CPO 内部的局部原子编号。数字 \verb|0| 表示空接口，不参与跨 CPO 连接。

\subsection{核心关键字}

核心片段由关键字序列构成。当前实现支持如下关键字：

\begin{table}[h]
\centering
\begin{tabular}{lll}
\toprule
类别 & 关键字 & 含义 \\
\midrule
单原子 & \verb|C P O N S F Cl Br I| & 对应元素原子 \\
常用取代基 & \verb|Me Ac NH2 OH SH OMe OPh =O| & 预定义官能团或片段 \\
常见环系 & \verb|Ph Im Ind Pyr| & 常用环片段 \\
参数化骨架 & \verb|nR nL| & $n$ 元碳环或 $n$ 原子碳链 \\
键阶修饰 & \verb|= #| & 双键或三键修饰符 \\
掺杂修饰 & \verb|^X| & X杂原子进入环系 \\
\bottomrule
\end{tabular}
\end{table}


\subsection{CPO 间接口：连边与稠合}

接口括号类型具有明确语义：
\begin{itemize}[leftmargin=2em]
  \item 圆括号 \verb|(p1,p2)| 表示直接连边接口。相邻 CPO 的对应接口原子按位置逐对相连。
  \item 花括号 \verb|{p1,p2}| 表示共用原子接口。相邻 CPO 的对应接口原子被视为同一组原子，用于稠合环。
\end{itemize}

例如：

\begin{lstlisting}[language=mychem]
(0)2L,(1,2);(1,2)2L,(0);
\end{lstlisting}

该表达式将第一个 CPO 的 1、2 号原子分别与第二个 CPO 的 1、2 号原子直接相连，形成两条跨 CPO 键。相反：

\begin{lstlisting}[language=mychem]
(0)6R,{4,5};{1,2}6R,(0);
\end{lstlisting}

表示第二个六元环的 1、2 号原子与第一个六元环的 4、5 号原子共用，从而构造稠合环。相邻接口必须使用相同括号类型，否则语义翻译阶段报错。旧规则中“两个非零接口自动稠合”的推断已被删除，因此 \verb|(4,5);(1,2)| 现在表示两条直接跨片段键，而不是稠合。

\subsection{取代基与原子替换}

CPO 的 \verb|SUBS| 部分描述当前核心上的局部修改：

\begin{lstlisting}[language=mychem]
p-X,          attach: 在 p 号原子外接 X
p1,p2-X,      attach many: 在多个位置外接 X
p-=,          升级 p 到 defaultNext[p] 的内部键为双键
p-=O,         在 p 号原子上外接双键氧
p-^X,         replace: 用 X 的入口原子替代 p 号原子
p-[CPO],      外接结构化 CPO 片段
p-^[CPO],     用结构化 CPO 的入口原子替代 p 号原子
\end{lstlisting}

\verb|p-X| 和 \verb|p-^X| 的区别是语言设计中的关键点。前者在原子 $p$ 外新增取代基，原子 $p$ 保持不变；后者将 $p$ 号原子本身替换为 $X$ 的入口原子，并保留 $p$ 原有的键。例如：

\begin{lstlisting}[language=mychem]
(0)6R,1-O,(0);     # 在 1 号碳上外接氧
(0)6R,1-^O,(0);    # 将 1 号碳替换为氧杂原子
\end{lstlisting}

结构化取代基使用方括号包围一个完整 CPO：

\begin{lstlisting}[language=mychem]
(0)6R,3-[(1)COH,(0)],(0);
\end{lstlisting}

组取代基要求非零左接口唯一，从而确定该组连接到母体的位置。

\section{词法与语法分析实现}

\subsection{词法分析}

词法分析器按最长匹配规则识别关键字。由于 \verb|OMe|、\verb|OPh|、\verb|OH| 与单原子 \verb|O| 共享前缀，关键字表将长关键字排在短关键字之前，避免将 \verb|OMe| 错误拆成 \verb|O| 与标识符。

当前 token 类型包括：
\begin{table}[h]
\centering
\begin{tabular}{lll}
\toprule
Token & 例子 & 用途 \\
\midrule
\verb|Sep| & \verb|,| & 分隔接口编号或取代项 \\
\verb|Semi| & \verb|;| & 结束一个 CPO \\
\verb|LParen/RParen| & \verb|(|, \verb|)| & 直接连边接口 \\
\verb|LBrace/RBrace| & \verb|{|, \verb|}| & 共用原子接口 \\
\verb|LStr/RStr| & \verb|[|, \verb|]| & 结构化取代基 \\
\verb|Keyword| & \verb|Ph|, \verb|=O| & 化学片段关键字 \\
\verb|Digit| & \verb|0|, \verb|12| & 局部位置编号 \\
\verb|Haf| & \verb|-| & 位置与取代基之间的连接标记 \\
\verb|Caret| & \verb|^| & 原子替换标记 \\
\bottomrule
\end{tabular}
\end{table}

词法阶段也承担部分早期错误检查。例如当读到数字后紧跟大写字母且不是 \verb|R| 或 \verb|L| 时，直接拒绝 \verb|5O| 这类旧补丁语法，并提示用户改用 \verb|5-^O|。

\subsection{递归下降 Parser}

Parser 使用手写递归下降实现。每个非终结符对应一个解析函数，例如 \verb|parseCpd|、\verb|parseCpo|、\verb|parseInterface|、\verb|parseSubs|、\verb|parseCore|。该选择的原因是：语言规模小、文法结构清晰，并且手写 parser 便于给出面向 DSL 用户的错误信息。

AST 的核心数据结构包括：
\begin{itemize}[leftmargin=2em]
  \item \verb|Cpd|：完整程序，由多个 CPO 组成；
  \item \verb|Cpo|：左接口、核心、取代基、右接口；
  \item \verb|Interface|：接口编号与接口类型 \verb|Connect/Fuse|；
  \item \verb|Core|：关键字序列；
  \item \verb|Subs|：取代基链表，记录 attach、replace、group 等信息。
\end{itemize}

语法分析只保证输入符合上下文无关结构。诸如接口编号是否存在、稠合原子类型是否匹配、端基是否继续延展等问题留给语义翻译阶段检查。

\section{IR 与语义设计}

\subsection{IR 节点}

IR 是显式无序连接图。每个节点表示一个原子：

\begin{lstlisting}
ATOM(self_id, conn1, conn2, conn3, conn4, conn5, conn6)
\end{lstlisting}

\verb|self_id| 是唯一原子编号；\verb|conn_i| 是相邻原子编号，\verb|0| 表示空连接槽。键阶通过重复邻接点表示：单键出现一次，双键出现两次，三键出现三次。例如 \verb|O(3,2,2,0,0,0,0)| 表示 3 号氧与 2 号原子形成双键。

\subsection{Fragment 属性}

AST 到 IR 的翻译以 \verb|Fragment| 为核心综合属性：

\begin{lstlisting}
Fragment {
  vector<IrNode> ir;
  int entry;
  optional<int> exit;
  int entryBond;
  vector<int> interfaceL, interfaceR;
  InterfaceKind interfaceLKind, interfaceRKind;
  map<int,int> defaultNext;
}
\end{lstlisting}

其中：
\begin{itemize}[leftmargin=2em]
  \item \verb|entry| 表示片段被外部接入时使用的入口；
  \item \verb|exit| 表示线性拼接时的出口，\verb|nullopt| 表示端基，不可继续作为 CORE 尾部延展；
  \item \verb|entryBond| 表示外部连接到入口时使用的键阶；
  \item \verb|interfaceL/interfaceR| 保存 CPO 左右接口编号；
  \item \verb|interfaceLKind/interfaceRKind| 保存接口语义，即直接连边或共用原子；
  \item \verb|defaultNext| 保存 \verb|p-=| 这类内部升键操作的默认目标。
\end{itemize}

\subsection{关键语义操作}

翻译器定义若干基本图操作：
\begin{itemize}[leftmargin=2em]
  \item \verb|shift(F, k)|：将片段中所有非零 id 和连接编号平移 $k$，避免合并时编号冲突；
  \item \verb|appendIr|：合并两个不冲突的 IR 节点集合；
  \item \verb|link(a,b,order)|：插入指定键阶的双向连接；
  \item \verb|ensureBond|：将已有键提升到指定键阶；
  \item \verb|rewrite(alias)|：在稠合或替换时重写原子编号。
\end{itemize}

这些操作共同保证 IR 始终保持双向连接一致性，并在连接槽不足时报告 valence overflow。

\section{AST 到 IR 的翻译}

\subsection{核心片段翻译}

每个关键字首先由 \verb|Keyword.cpp| 翻译为基本片段。例如：
\begin{itemize}[leftmargin=2em]
  \item \verb|C|、\verb|O|、\verb|N| 等单原子关键字翻译为单节点片段；
  \item \verb|Me|、\verb|OH|、\verb|NH2| 翻译为预定义官能团；
  \item \verb|nR| 翻译为 $n$ 元环，并建立环上 \verb|defaultNext|；
  \item \verb|nL| 翻译为 $n$ 原子链；
  \item \verb|=| 和 \verb|#| 不产生独立节点，而是作为后续连接的键阶修饰。
\end{itemize}

当 CORE 包含多个关键字时，翻译器从左到右合并片段。若左侧片段 \verb|exit = nullopt|，说明该片段是端基，不能继续作为线性 CORE 的左侧延展，此时翻译器报错。

\subsection{取代基翻译}

对于普通取代 \verb|p-X|，翻译器执行：
\begin{enumerate}[leftmargin=2em]
  \item 翻译 $X$ 为片段 $S$；
  \item 将 $S$ 的编号平移到不与母体冲突；
  \item 找到 $S$ 的入口；
  \item 将母体 $p$ 与入口按 \verb|entryBond| 相连；
  \item 合并 $S$ 的 IR 和 \verb|defaultNext|。
\end{enumerate}

对于原子替换 \verb|p-^X|，翻译器不新增一条外部键，而是把 $X$ 的入口原子映射到母体 $p$：
\begin{enumerate}[leftmargin=2em]
  \item 要求 $p$ 在母体中存在，且 $X$ 的 IR 非空；
  \item 翻译并平移 $X$；
  \item 选取 $X$ 的入口原子；
  \item 将母体 $p$ 的元素类型替换为入口原子的元素类型；
  \item 保留 $p$ 的原有连接；
  \item 将 $X$ 中除入口外的节点和内部键合并到母体。
\end{enumerate}

该设计使氧杂、氮杂、硫杂等环内原子替换可以统一写成 \verb|p-^O|、\verb|p-^N|、\verb|p-^S|，避免引入 \verb|5O| 这类临时补丁语法。

\subsection{CPO 间连接与稠合}

相邻 CPO 的组合由左片段右接口和右片段左接口决定。若二者括号类型不同，翻译器报错。

\paragraph{直接连边。}
当接口类型为 \verb|Connect| 时，要求两侧接口长度相等且不含 0。翻译器平移右片段编号，合并 IR，并对每一对对应接口原子调用 \verb|link|：
\[
  \forall i,\quad
  \operatorname{link}(A.\mathrm{interfaceR}_i, B.\mathrm{interfaceL}_i, B.\mathrm{entryBond}).
\]

\paragraph{共用原子稠合。}
当接口类型为 \verb|Fuse| 时，要求两侧各有两个非零接口。翻译器构造别名映射：
\[
  B.\mathrm{interfaceL}_1 \mapsto A.\mathrm{interfaceR}_1,\quad
  B.\mathrm{interfaceL}_2 \mapsto A.\mathrm{interfaceR}_2.
\]
随后将右片段所有节点与边通过该映射重写并合并。若同一个 IR id 对应的原子类型不一致，例如一个是 \verb|N| 而另一个是 \verb|C|，则报告 ``Fuse atom mismatch''。这可以阻止不合理的稠合表达式静默生成错误结构。

\section{渲染后端}

\subsection{IR 到分子对象}

渲染模块首先从 IR 中恢复边集合，并根据重复连接次数计算键阶。若启用 RDKit，系统将 IR 降低为 \verb|RDKit::RWMol|：
\begin{enumerate}[leftmargin=2em]
  \item 根据 IR 节点元素类型创建 RDKit Atom；
  \item 根据 IR 边集合创建 RDKit Bond；
  \item 调用 \verb|updatePropertyCache(false)|；
  \item 调用 RDKit 2D 坐标生成；
  \item 使用 \verb|MolDraw2DSVG| 输出 SVG。
\end{enumerate}

RDKit 后端的优点是布局质量高、能处理常见化学绘图约定；缺点是依赖外部库。项目的 CMake 配置中对 RDKit 采用手动库发现，以规避部分系统包导出的 CMake target 引用缺失 NumPy 或 Cairo 路径的问题。由于化学物质的极度多样性，后端需要处理大量实际上与编译无关的特例，所以这里引入避免重复工作。

\section{使用方式}


项目将可复用逻辑封装在 \verb|compiler_core| 静态库中。其上有多个入口：
\begin{itemize}[leftmargin=2em]
  \item \verb|tokenizer|：命令行词法调试工具；
  \item \verb|parser|：命令行完整编译与渲染工具，会打印 token、AST、IR 和 SVG；
  \item \verb|chem_http_server|：监听本地端口 \verb|127.0.0.1:17654|，提供 \verb|POST /render|；
  \item \verb|chem_native_host|：Chrome/Edge native messaging 备用入口。
\end{itemize}

浏览器插件属于前端外壳。其任务是在网页 DOM 中识别 \verb|<my-Chem>| 标签或被 HTML 转义显示的 my-Chem 文本，提取 DSL 内容，通过 HTTP 调用本地 C++ 服务，并用返回的 SVG 替换原位置。插件侧不实现化学语义，避免前后端重复逻辑；所有语法和语义错误均由后端统一返回。对于单个标签解析失败的情况，插件应跳过该标签并继续处理其它标签，以保证页面鲁棒性。\textbf{后两部分与编译无关，且受限于个人前端水平，实现效果不佳}

\section{表达能力与Demo}
\begin{itemize}


\item 任何物质通过足够多次的切断都变成链状骨架进而被表达，但是只是直观性可能会损失。
\item 大模型具有一定生成此DSL的能力，只通过对话内的提示（提示词见仓库）加上一些迭代即可完成较复杂的绘图工作，如果使用skill可能可以达到更好的效果。\url{https://chatgpt.com/share/6a46a5a1-6da0-83e8-98e6-7df919d9aca1}。
\item 对于简单的分子，大模型很容易写出DSL，例如对于所有构成蛋白质的氨基酸：
\end{itemize}

\subsection{简单分子示例}

下面的 SVG 图由 \verb|test_cases_amino_acids.tsv| 渲染得到。所有图片均使用仓库相对路径 \verb|output/|，避免依赖本机绝对路径。

\begin{figure}[htbp]
\centering
\begin{tabular}{ccccc}
\includesvg[width=0.16\textwidth]{output/G_Gly} &
\includesvg[width=0.16\textwidth]{output/A_Ala} &
\includesvg[width=0.16\textwidth]{output/V_Val} &
\includesvg[width=0.16\textwidth]{output/L_Leu} &
\includesvg[width=0.16\textwidth]{output/I_Ile} \\
Gly & Ala & Val & Leu & Ile \\
\includesvg[width=0.16\textwidth]{output/S_Ser} &
\includesvg[width=0.16\textwidth]{output/T_Thr} &
\includesvg[width=0.16\textwidth]{output/C_Cys} &
\includesvg[width=0.16\textwidth]{output/M_Met} &
\includesvg[width=0.16\textwidth]{output/D_Asp} \\
Ser & Thr & Cys & Met & Asp \\
\includesvg[width=0.16\textwidth]{output/E_Glu} &
\includesvg[width=0.16\textwidth]{output/N_Asn} &
\includesvg[width=0.16\textwidth]{output/Q_Gln} &
\includesvg[width=0.16\textwidth]{output/K_Lys} &
\includesvg[width=0.16\textwidth]{output/R_Arg} \\
Glu & Asn & Gln & Lys & Arg \\
\includesvg[width=0.16\textwidth]{output/F_Phe} &
\includesvg[width=0.16\textwidth]{output/Y_Tyr} &
\includesvg[width=0.16\textwidth]{output/W_Trp} &
\includesvg[width=0.16\textwidth]{output/H_His} &
\includesvg[width=0.16\textwidth]{output/P_Pro} \\
Phe & Tyr & Trp & His & Pro \\
\end{tabular}
\caption{MarkChem 渲染的氨基酸简单示例。}
\end{figure}

\subsection{复杂分子示例}

下面的 SVG 图由 \verb|test_cases_complex.tsv| 渲染得到。所有图片均使用仓库相对路径 \verb|output_complex/|。

\begin{figure}[htbp]
\centering
\begin{tabular}{cccc}
\includesvg[width=0.22\textwidth]{output_complex/VitC_Vitamin_C} &
\includesvg[width=0.22\textwidth]{output_complex/VitD_Vitamin_D} &
\includesvg[width=0.22\textwidth]{output_complex/B12_Vitamin_B12_corrin_like} &
\includesvg[width=0.22\textwidth]{output_complex/Car_Beta_carotene} \\
维生素 C & 维生素 D & 维生素 B12 近似 & 胡萝卜素 \\
\includesvg[width=0.22\textwidth]{output_complex/Chl_Chlorophyll_no_Mg} &
\includesvg[width=0.22\textwidth]{output_complex/Lut_Lutein} &
\includesvg[width=0.22\textwidth]{output_complex/CTX_Ciguatoxin_ladder_polyether_like} &
\includesvg[width=0.22\textwidth]{output_complex/Taxol_Paclitaxel_taxane_like} \\
叶绿素近似 & 叶黄素 & 雪卡毒素近似 & 紫杉醇近似 \\
\includesvg[width=0.22\textwidth]{output_complex/EE2_Ethinylestradiol_steroid_like} &
\includesvg[width=0.22\textwidth]{output_complex/Colch_Colchicine_tropolone_like} &
\includesvg[width=0.22\textwidth]{output_complex/Norb_Norbornene} &
\includesvg[width=0.22\textwidth]{output_complex/Hel6_Hexahelicene_like} \\
乙炔雌二醇近似 & 秋水仙素近似 & 降冰片烯 & 六螺苯近似 \\
\includesvg[width=0.22\textwidth]{output_complex/Met_Metformin} &
\includesvg[width=0.22\textwidth]{output_complex/PenG_Penicillin_G_like} &
 &
 \\
二甲双胍 & 青霉素 G 近似 & & \\
\end{tabular}
\caption{MarkChem 渲染的复杂分子示例。}
\end{figure}
\section{代码仓库与环境配置}
代码位于仓库 \url{https://github.com/lambda-z-miu/compiler_project} 
需配置RDKit环境，推荐在Linux下进行，作为参考本地开发时使用的WSL Ubuntu 24.

\end{document}
